
%\needspace{4\baselineskip}
\color{uniblue}{\subsection{Контроль отсутствия напряжения на 1(2) секции шин (КОН 1сш, КОН 2сш)}}
\color{black}

\begin{enumerate}[label=\arabic{section}.\arabic{subsection}.\arabic*, labelsep=4pt, leftmargin=0pt, itemindent=57pt]

\item
Функция контролирует отсутствие напряжения на 1(2) секции шин. В устройстве контроль отсутствия напряжения на 1(2) секции шин представлен функциями «КОН 1сш» и «КОН 2сш» соответственно. Функции используется для АПВ, а также для включения с контролем синхронизма.

\item
Пусковые условия функций формируются без выдержки времени при снижении всех линейных напряжений ниже уставки и одновременном выполнении следующих условий если уставка SGF2 выставлена в положение «Измеренный»:
\begin{itemize}
    \item функция введена в работу;
    \item минимальное из подведенных линейных напряжений, измеренных с 1(2) секции шин ниже уставки;
    \item отсутствует срабатывание ИО напряжения обратной последовательности;
    \item не активный сигнал срабатывания БНН 1(2)cш.
\end{itemize}

\item
Предусмотрена блокировка функции при срабатывании ИО обратной последовательности напряжения выше уставки и наличии сигнала срабатывания БНН 1(2) секции шин.

\item
Если уставка SGF2 выставлена в положение «Внешний» работа функции осуществляется при наличии активного внешнего дискретного сигнала «Наличие U 1(2)сш» без учета срабатывания измерительных органов.

\item
Ввод функции в работу на этапе параметрирования устройства осуществляется программным переключателем SGF1 «Ввод\_функции». Информация о введенной функции отображается выходным сигналом «КОН 1(2)сш: Ввод».

\item
Функция «КОН 1(2)сш» может быть выведена из работы оперативно путем активации сигнала «ОВ КОН 1(2)сш», что характеризуется активным состоянием сигнала «КОН 1(2)сш: Оперативный вывод» на выходе функции. 

\item
Алгоритм работы функций выполнен в соответствии с рисунком \ref{konsh:img1}. В таблице \ref{konsh:tbl1} приведены параметры, необходимые для настройки КОН 1сш. В таблице \ref{konsh1:tbl1} приведены параметры, необходимые для настройки КОН 2сш.

\vspace{3mm}
\begin{figure}[H]
\centering
\includegraphics[width=0.85\textwidth,height=0.85\textheight,keepaspectratio]{\fbpath/408. КОН сш/_latex/img/konsh.pdf}
\captionof{figure}{Функциональная схема логики функции «КОН 1(2)сш»}\label{konsh:img1}
\end{figure}

\end{enumerate}
